\chapter{Literature Review}

\section{Existing Approaches in Practice and Academia}

\(z\)-buffering is a contemporary tool for approximating occlusion 
of triangles by taking a pixel-level raster of rays from the 
viewing plane onto what lies before it. The colour of the first 
triangle in the path of such a ray is assigned to the pixel of the 
ray. There is also the painter's algorithm in which primitives are 
layered from back to front, though it is a very general description
\cite{nyberg2011hidden}. One variant of the painter's algorithm 
in the context of parametric renderings is the centroid sort, in 
which geometric primitives are ordered by the depth of the
average of their vertices \cite{skillmon2025answer3dpointsort}.
There are two ways to find these depths; one way is to take the 
dot product with the observer vector \cite{5062592}. Another way 
is to use an affine transformation to align the \(z\) axis with 
the observer vector, and then order by just that one component.

Setting the camera position aand orientation is done by determining
the affine transformation matrix which carries the camera to the 
desired position, and performing the inverse of that transformation 
on all of the parametric objects. Perspective is achieved by using 
a projective transformation after all the affine transformations 
are done. While this practise is commonplace in 3D graphics, it has 
yet to be achieved using tools in the \TeX{} ecosystem.

\section{Limitations and Open Challenges in Current Methods}

\(z\)-buffering is fundamentally an approximation. While it is 
useful for real-time 3D graphics, the author of this book contends
that aa definitive occlusion sorting of non-intersecting and 
non-cyclically overlapping primitives should be considered as a 
fundamental concept in 3D occlusion. The centroid sort on the other 
hand is prone to non-subtle occlusion errors 
\cite{jasper2025tikzparametricsurface}. In attempt to resolve this,
I tried separating the triangles into groups based on the orientation
of their normal vector with respect to the observer's direction, 
and then performed a centroid sort on both groups. This method was 
still insufficient.

\section{Proposed Approach and Its Advantages Over Existing Work}

This package proposes a transitive occlusion comparator for 
non-intersecting and non-cyclically overlapping geometric primitives.
This comparator is designed to be used alongside an automatic 
primitive-primitive clipping algorithm, and with a more sophisticated 
variant of back-face culling which accounts for translucency.
The future goal of this package is to introduce the aforementioned 
geometric primitive clipping software to complement this sorting 
algorithm. Possibly, another goal is to implement such a back-face 
culling algorithm as well, though that is secondary to the current 
goal. The advantage of the occlusion comparator introduced in this 
package is its clarity. Unlike \(z\)-buffering, this method offers 
a definitive ordering of non-intersecting and non-cyclically 
overlapping geometric primitives. It is paired with an artificial
inteligence-written topological sort function which is capable of 
ordering the primitives based on the comparator's output.